\HMTNarrativeHeading{Lo que conserva una transformación}

La teoría operatoria da una forma calculable a la conservación del
estado y de sus registros. Una lectura parcial obtiene una componente;
el complemento recoge la contribución que permite restituir el
balance completo. La suma de las normas cuadradas de la componente leída
y de su complemento reproduce la norma cuadrada del estado inicial.
Las fórmulas de recuperación especifican qué
datos y qué operadores hacen posible invertir la lectura.

Al iterar el transporte, cada etapa incorpora un complemento y
prolonga una componente terminal. El balance finito conserva
explícitamente ese terminal. El paso a profundidad arbitraria
estudia la convergencia y determina cuándo el archivo reúne toda la
norma. Este orden permite interpretar la memoria como una operación
de conservación y recuperación, con condiciones verificables.

La extensión exterior muestra cómo se transportan áreas y volúmenes
de la realización lineal. Las contribuciones mixtas participan en
los productos exteriores y en sus balances. El cambio de
representación conserva las relaciones que esos objetos miden.
La recuperación mantiene el enlace con la orientación y las formas
que utiliza cada prueba.

La monodromía y los lazos entre regímenes geométricos añaden una
dinámica determinada. El artículo compone transformaciones elípticas,
parabólicas e hiperbólicas, conserva su memoria y calcula su
respuesta. Las simetrías de la acción de transporte producen las
cargas variacionales de ese dominio. La invariancia de un observable
se formula, por su parte, mediante su relación con el transporte.
Estas dos formas de conservación se enlazan conservando sus
definiciones particulares.

Las realizaciones físicas recuperan después las secciones de acción,
los lectores de área y las escalas previamente generadas. El ejemplo
gaussiano permite calcular la pureza y el espectro de una lectura
reducida desde una preparación conjunta especificada. Sus
correlaciones conservan la información del estado completo. Este
resultado muestra una consecuencia operativa del argumento: el
transporte determina tanto lo que una observación recibe como la
información necesaria para comprender su relación con la totalidad.
